\documentclass[../main.tex]{subfiles}
\ifSubfilesClassLoaded{\setcounter{chapter}{3}}{}
\begin{document}

\chapter{Missing Values, Duplikat, dan Outlier Dasar}

\begin{subcpmk}
  \item \textbf{Sub-CPMK-P3:} Menangani missing, duplikat, outlier sederhana dengan narasi keputusan (RPS Mg.\ 3).
\end{subcpmk}

\SesiRPS{3}{P3}{$\sim$170 menit}{CPMK-P2}

\section{Kemampuan akhir sesi}

Melaporkan NA per kolom; memilih strategi imputasi atau penghapusan; menghapus duplikat; mengidentifikasi outlier visual (boxplot) dan mendokumentasikan tindakan.

\section{Teori}

\begin{konsep}
\textit{Missing} dapat MCAR/MAR/MNAR---pada praktikum, yang penting Anda \textbf{transparan}: berapa banyak yang hilang, apa yang Anda lakukan, dan risiko bias. Duplikat baris dapat merusak statistik dan model; outlier bukan selalu error---bisa sinyal---tetapi harus diputuskan dengan konteks domain.
\end{konsep}

\section{K3 / etika}

Jangan menghapus baris berisi kelompok minoritas tanpa analisis dampak terhadap keadilan model (catatan untuk proyek akhir).

\section{Sarana}

\texttt{pandas}, \texttt{numpy}, \texttt{matplotlib.pyplot} atau \texttt{seaborn} untuk boxplot.

\section{Prosedur kerja}

\begin{enumerate}[leftmargin=*]
  \item Hitung \texttt{isna().sum()} per kolom; simpan salinan \texttt{df\_raw}.
  \item Terapkan strategi (imputasi median/mode atau \texttt{dropna} selektif) dengan komentar Markdown.
  \item \texttt{drop\_duplicates()}; laporkan berapa baris yang hilang.
  \item Plot boxplot satu kolom numerik; tulis keputusan outlier.
  \item Simpan \texttt{df\_clean} untuk minggu berikutnya.
\end{enumerate}

\section{Contoh kode}

\begin{lstlisting}[caption=Audit missing values]
import pandas as pd
import numpy as np

# df = pd.read_csv("...")
missing = df.isna().sum().sort_values(ascending=False)
print(missing[missing > 0])
\end{lstlisting}

\begin{lstlisting}[caption=Imputasi median untuk kolom numerik]
num_cols = df.select_dtypes(include="number").columns
df_clean = df.copy()
for c in num_cols:
    df_clean[c] = df_clean[c].fillna(df_clean[c].median())
\end{lstlisting}

\begin{lstlisting}[caption=Duplikat dan shape]
print("sebelum", df_clean.shape)
df_clean = df_clean.drop_duplicates()
print("sesudah", df_clean.shape)
\end{lstlisting}

\begin{lstlisting}[caption=Boxplot outlier (contoh)]
import matplotlib.pyplot as plt
col = df_clean.select_dtypes(include="number").columns[0]
plt.boxplot(df_clean[col].dropna())
plt.title(f"Boxplot: {col}")
plt.show()
\end{lstlisting}

\section{Referensi}

\begin{itemize}[leftmargin=*]
  \item pandas missing data. \url{https://pandas.pydata.org/docs/user_guide/missing_data.html}
  \item Kaggle Intermediate ML. \url{https://www.kaggle.com/learn/intermediate-machine-learning}
\end{itemize}

\begin{aktivitas}
  \item Bandingkan \texttt{shape} sebelum/sesudah; tulis paragraf \textit{trade-off}.
\end{aktivitas}

\begin{latihan}
  \item Kapan median lebih aman daripada mean untuk imputasi?
\end{latihan}

\begin{asesmen}
\textbf{Tugas modul:} \texttt{df\_clean} + narasi keputusan; bukti visual outlier minimal satu kolom.
\end{asesmen}

\begin{checklist}
  \item Saya menghitung dan menangani missing secara dokumentatif
  \item Saya memeriksa duplikat
\end{checklist}

\begin{rangkuman}
Minggu 3 membersihkan data secara terdokumentasi untuk mencegah ``angka ajaib'' tanpa alasan.

\textbf{Kata kunci:} \texttt{fillna}, \texttt{dropna}, \texttt{drop\_duplicates}, boxplot
\end{rangkuman}

\end{document}
